% 19 November 2026. Drafted for Edward Baker with OpenAI Codex assistance.
% Lead actual backend/effort unavailable; helpers explicitly Sol6.1/ultra.
% Analytical integration accepted in Review05; numerical integration assessed separately.
\section{Retaining the arithmetic correction on the fixed cell}
\label{sec:corrected-reference}

The correction in Section~6 changes the approximation center. We now
make that change on the same cell as \eqref{le:pilot-cell}, paying the
difference between the published value formula and a holomorphic reference.
This is a source-acquisition refinement. Direct evaluation of the corrected
reference receives the same information and complete allowance.

\subsection{A local holomorphic reference and its derivative bounds}

Put $k=\pi/8$, $w=z/2+kt$ and
\begin{align}
 p(w)&=1-2\bigl(\sqrt{w/(2\pi)}-N\bigr),\notag\\
 C_0(p)&=\frac{e^{\pi\ii(p^2/2+3/8)}-\ii\sqrt2\cos(\pi p/2)}
                   {2\cos(\pi p)},\qquad
 C_0^\sharp(p)=\overline{C_0(\bar p)},\notag\\
 A_N(w)&=\frac{\sqrt w}{2}
       \bigl\{(w+\ii)C_0(p(w))+(w-\ii)C_0^\sharp(p(w))\bigr\},\notag\\
 J_N(t,z)&=(-1)^{N+1}\frac{\sqrt\pi}{4}e^{-kz-k^2t},\qquad
 C_{\rm hol}(t,z)=J_N(t,z)A_N(w),\qquad Q_N=P_N-C_{\rm hol}.
 \label{cr:reference}
\end{align}
The square roots are principal. All assertions in this section concern
$N=33$, $|\Re z-14010.1|\le2$, $|\Im z|\le1$ and
$9/100\le t\le1/10$; no global continuation of the displayed quotient
is needed.

\begin{proposition}[Explicit correction bounds]
\label{cr:bounds}
On this rectangle the reference is holomorphic in $z$ and respects
conjugation. Its slow profile satisfies
\begin{equation}
 |A_N|<1060000,\qquad |A_N'|<24000,\qquad |A_N''|<1400.
 \label{cr:profile-bounds}
\end{equation}
\end{proposition}

\begin{proof}
The rational bounds $157/50<\pi<22/7$ give
$7004<\Re w<7006.09$, $|\Im w|\le1/2$ and $|w|<7007$.
Writing $\sqrt{w/(2\pi)}=a+\ii b$ and using
$a^2=(|w|+\Re w)/(4\pi)$ and $2ab=\Im w/(2\pi)$ gives
\begin{equation}
 \frac9{50}<\Re p<\frac6{25},\qquad |\Im p|<\frac1{400},
 \qquad |p|<\frac14,\qquad |\cos(\pi p)|>\frac7{10}.
 \label{cr:p-domain}
\end{equation}
The last inequality follows from
$|\cos(u+\ii v)|^2=\cos^2u+\sinh^2v$ and
$\cos(6\pi/25)>7/10$. Thus all denominators and branches remain
clear on a neighborhood of the closed rectangle.

For completeness, set $n=e^{\pi\ii(p^2/2+3/8)}-
\ii\sqrt2\cos(\pi p/2)$ and $d=2\cos(\pi p)$.
On \eqref{cr:p-domain}, the exponential has modulus less than
$1003/1000$ and $\cosh(\pi|\Im p|)<1001/1000$. Using
$|\sin z|\le|z|\cosh|\Im z|$ and $\sqrt2<10/7$ gives
$|n|<5/2$, $|n'|<7/4$, $|n''|<15/2$,
$|d'|<5$, $|d''|<20$ and $|d|>7/5$. Thus
$|C_0|<9/5$, $|C_0'|<8$ and $|C_0''|<90$;
the quotient identities are
\[
 C_0'=\frac{n'-d'C_0}{d},\qquad
 C_0''=\frac{n''-2d'C_0'-d''C_0}{d},\qquad d''=-\pi^2d.
\]
The same bounds hold for $C_0^\sharp$. Write $A_N=gB$, where
$g=\sqrt w/2$ and $B$ is the expression in braces in
\eqref{cr:reference}. Then
\[
 |g|<42,\quad |g'|<\frac1{4\cdot83},\quad
 |g''|<\frac1{8\cdot83^3},\quad
 |p'|<\frac1{200},\quad |p''|<\frac1{2800000}.
\]
Since $|w\pm\ii|<7008$, product and chain rules give
\begin{align*}
 |B|&<2\cdot7008\cdot\frac95,\\
 |B'|&<\frac{18}5+2\cdot7008\cdot\frac8{200},\\
 |B''|&<\frac{32}{200}
   +2\cdot7008\left(\frac{90}{200^2}
                        +\frac8{2800000}\right).
\end{align*}
Substitution into $A_N'=g'B+gB'$ and
$A_N''=g''B+2g'B'+gB''$ proves \eqref{cr:profile-bounds}.
Finally $p(\bar w)=\overline{p(w)}$ and the paired sharp terms give
$A_N(\bar w)=\overline{A_N(w)}$. This also proves the claimed
physical conjugation symmetry.
\end{proof}

Let $T'=x/2+kt$ and let $C_{\rm val}$ denote the correction in
Corollary~6.4 of \cite{Polymath2019}. Substituting the principal-branch
formula for $\mathfrak M_0$ and the source's phase
$U(T')=\exp[-\ii\{(T'/2)\log(T'/(2\pi))-T'/2-\pi/8\}]$
gives exactly
\begin{equation}
 \mathfrak M_0(\ii T')U(T')e^{\pi\ii/8}
 =-\frac{\sqrt\pi}{8}e^{-\pi T'/4}\sqrt{T'}(T'+\ii),
 \qquad C_{\rm val}(t,x+\ii y)=J_N(t,x+\ii y)A_N(T').
 \label{cr:phase}
\end{equation}
Consequently the real-axis corrections agree. Off-axis their difference
is the vertical integral of $A_N'$, which is paid below. Differentiating
the holomorphic reference, with $w_t=k$, $w_z=1/2$, gives the exact
cancellations
\begin{equation}
 (\partial_t+\partial_z^2)C_{\rm hol}
       =\frac14J_NA_N''(w),\qquad
 \mathcal R_Q=\mathcal R_N-\frac14J_NA_N''(w).
 \label{cr:defect}
\end{equation}
The $A_N$ and $A_N'$ terms cancel. In backward time the corrected
error satisfies $(H-Q)_u=(H-Q)_{zz}+\mathcal R_Q$.

\subsection{The complete allowance and its final consumer}

Define the positive leading payment
\begin{equation}
 L(t,x)=e^{k^2t}|\mathfrak M_0(\ii T')|
       =\frac{\sqrt\pi}{8}e^{-kx-k^2t}
                              \sqrt{T'}\sqrt{T'^2+1}.
 \label{cr:leading}
\end{equation}
Write $\mathcal V_P$ for \eqref{le:source-value}. Corollary~6.4,
together with the vertical mismatch, supplies the complete bound
\begin{equation}
 |H-Q_N|\le\mathcal V_Q
 :=\mathcal V_P-L+\Delta,\qquad
 \Delta=|J_N(t,z)|\frac{|y|}{2}\,24000.
 \label{cr:source}
\end{equation}
Both summand errors and the entire retained $\delta_-+\delta_+$
remainder remain in $\mathcal V_P-L$. Equations
\eqref{cr:profile-bounds} and \eqref{cr:leading} give
\begin{equation}
 \frac{\Delta}{L}<\frac{24000}{83\cdot7004}<\frac1{24},
 \qquad \frac{|C_{\rm hol}|}{L}
       <\frac{2\cdot1060000}{83\cdot7004}<4.
 \label{cr:saving}
\end{equation}
Thus more than $23/24$ of the \emph{leading omitted payment} is saved.
This fraction is not a percentage improvement of the complete gap bound.

There is also a concrete count transfer. The previous certificate has
$\mathcal V_P/|P_N|<1/6$ on the circle. Hence
$|C_{\rm hol}/P_N|<2/3$, and Rouch\'e transfers the count two to $Q_N$.
Pointwise, putting $q_P=\mathcal V_P/|P_N|$, one obtains
\begin{equation}
 \frac{\mathcal V_Q}{|Q_N|}
 \le\frac{\mathcal V_P-(23/24)L}{|P_N|-4L}<q_P,
 \quad\text{since }4q_P<23/24.
 \label{cr:relative}
\end{equation}
These inequalities do not supply $Q_N$'s center or discriminant. A final
comparison must acquire them and evaluate
$D_Q-r^2F(\operatorname{arctanh}q_Q,|c_Q-X|/r)$ with every numerical
payment. No $P_N$ denominator or center is reused as though it belonged
to $Q_N$. Direct $Q_N$ evaluation receives the same source information.

The collision controls remain consistent. For an exact colliding
reference the error may vanish while its discriminant reaches zero;
the strict final margin then fails. Modified theta coefficients cannot
import \eqref{cr:source} unchanged, and a changed observation kernel
fails the physical heat equation. These premise failures do not assert
that every collision-free box of a control must fail. The present local
branch and cutoff conditions do not establish an all-height family,
a passage through cutoff changes, or new canonical coverage.

\subsection{The complete corrected comparison}

The second calculation keeps exactly the previous $192$-bit arithmetic,
$64$ closed time cells, $512$ closed arcs and three radii. It evaluates
$Q_{33}$ directly and reacquires its count, denominator and annular
logarithm from $\Re(Q_{33}(z)/(z-X)^2)>0$ on every boundary arc.
The two logarithmic moments use the same fully paid Laurent alias formula
\eqref{le:aliases}; interval time variation is included. Both term errors,
the retained remainder and $12000|J_N||y|$ enter the local-arc ratio.
The old $P_{33}$ clearance and moments are not substituted for $Q_{33}$'s.

\begin{center}
\begin{tabular}{lr}
\toprule
Quantity over the entire fixed interval & Outward summary\\
\midrule
Corrected relative-error allowance $q_Q$ & $<0.015234762$\\
Corrected reference discriminant $D_Q$ & $>0.1617006119$\\
Refined corrected transfer allowance & $<0.012369746$\\
E10 complete lower bound at $Q$ & $>0.1491891639$\\
Refined complete lower bound at $Q$ & $>0.1493308659$\\
Previous refined complete lower bound at $P$ & $>0.0054147305$\\
\bottomrule
\end{tabular}
\end{center}

The refined corrected lower bound exceeds the previous evaluated lower
bound in every one of the $64$ closed time cells. In particular
\begin{equation}
 D_H(t)>0.1493308659\qquad(9/100\le t\le1/10)
 \label{cr:pilot-margin}
\end{equation}
on the same counted circle. The real pair has gap greater than $0.7728$,
since $(0.7728)^2/4=0.14930496<0.1493308659$.
This is a stronger sufficient lower bound for the same canonical pair;
the zeros have not changed. The gain uses retained correction information,
and is shared by direct evaluation with the same information. It does not
show an advantage of a time-integrated representation.

The scaled complete source allowance is below $21410.476$ and the
corrected denominator exceeds $941112.53$. Their independent ratio is
not used in place of the sharper local-arc ratios. The complete corrected
defect is conservatively bounded by $|\mathcal R_N|+350|J_N|$;
its scaled contour bound is below $357.797450$. This is an upper
triangle estimate, not an evaluation of the signed corrected defect.
No second scalar heat-propagation comparison is claimed.
The code, full record, separate assessment and sole scratch replay are
saved in the investigation. They retain the imported source theorem and
internal LLM-review limits. The fixed cell remains inside known coverage;
no new height/time region or improved de Bruijn--Newman bound follows.
